% TOP_PROBLEM: {"id": 1283, "title": "Transcendence of Apéry's Constant", "category_id": 1, "category": {"id": 1, "name": "number_theory", "display_name": "Number Theory", "description": "Properties of integers, prime numbers, Diophantine equations.", "slug": "number-theory", "order_index": 1, "created_at": "2026-07-31T15:26:25.670Z"}, "status": "open"}
% FIRST_POSTED: 2026-09-25
% ENTRY_KIND: statement_only
% !TeX program = lualatex
% Definition and sources only; this file is not an LLM solution attempt.
\documentclass[11pt]{article}
\usepackage[margin=25mm]{geometry}
\usepackage{fontspec}
\setmainfont{Latin Modern Roman}
\usepackage{amsmath,amssymb,mathtools}
\usepackage{unicode-math}
\setmathfont{Latin Modern Math}
\usepackage{xurl,hyperref}
\hypersetup{colorlinks=true,urlcolor=blue}
\setcounter{secnumdepth}{0}
\setlength{\parindent}{0pt}
\setlength{\parskip}{5pt}
\setlength{\emergencystretch}{3em}
\begin{document}
\section*{\texorpdfstring{Transcendence of Apéry's Constant}{Transcendence of Apéry's Constant}}\label{1283}
\subsection{Definitions and mathematical statement}
Ap\'ery's constant is the absolutely convergent series
\[
 \zeta(3)=\sum_{n=1}^{\infty}\frac1{n^3}.
\]
A real or complex number $z$ is algebraic over ${\mathbb{Q}}$ if some nonzero polynomial in ${\mathbb{Q}}[T]$ vanishes at $z$, and transcendental otherwise. Thus the selected assertion is
$P(\zeta(3))\ne0$ for every $P\in{\mathbb{Q}}[T]\setminus\{0\}$.
A negative resolution would give algebraicity, not necessarily rationality. Irrationality excludes only polynomial relations of degree one; it does not by itself settle this transcendence question. The target concerns this single value, not algebraic independence of a collection of odd zeta values.
\par\smallskip\textit{Formulation and concept references:} \href{https://arxiv.org/abs/math/0206176}{[S1]}.

\subsection{Short English statement}
Is the sum of the reciprocal cubes of the positive integers transcendental, rather than merely irrational?
\subsection{Sources}
\begin{itemize}
\item[S1]W. Zudilin, Arithmetic of linear forms involving odd zeta values, J. Theor. Nombres Bordeaux 16 (2004).
\url{https://arxiv.org/abs/math/0206176}.
\textit{Formulation source.}
\end{itemize}
\end{document}
